\SeventhArticle{18}
Die Untersuchungen des vorigen Artikels zeigen zwar, wie sich geeignete Funktionen $U$ für die einzelnen Werte der Zahl $n$ finden lassen, jedoch nur der Reihe nach, da man von kleineren Werten zu größeren übergehen muss. Wir sehen aber leicht, dass diese Funktionen allgemein dargestellt werden durch
\[
\begin{aligned}&
u^{n+1}-\frac{(n+1)n}{2\mathbin{.}(2n+1)}u^{n-1}
\\[5pt]&+\frac{(n+1)n(n-1)(n-2)}{2\mathbin{.}4(2n+1)(2n-1)}u^{n-3}
\\[5pt]&-\frac{(n+1)n(n-1)(n-2)(n-3)(n-4)}{2\mathbin{.}4\mathbin{.}6\mathbin{.}(2n+1)(2n-1)(2n-3)}u^{n-5}\\ &
+\text{usw.}
\end{aligned}
\]
oder auch, wenn wir das Zeichen $F$ wie in der oben angeführten Abhandlung verwenden, durch
\[
u^{n+1}F(-\tfrac12n,\ -\tfrac12(n+1),\ -(n+\tfrac12),\ u^{-2})
\]
Diese Induktion lässt sich durch die allgemein bekannte Methode leicht in einen strengen Beweis verwandeln oder, wenn man es vorzieht, mithilfe der Formel 19 in der angeführten Abhandlung. Die Funktion $U$ kann, wenn man es vorzieht, auch bei umgekehrter Reihenfolge der Glieder dargestellt werden durch
\[
\pm\frac{3\mathbin{.}5\mathbin{.}7\mathinner{.\,.\,.\,.}(n+1)}{(n+3)(n+5)\mathinner{.\,.\,.\,.}(2n+1)}\mathbin{.}uF(-\tfrac12n,\ \tfrac12(n+3),\ \tfrac32,\ uu)
\]
für gerades $n$, wobei das obere oder untere Zeichen gilt, je nachdem, ob $\frac12n$ gerade oder ungerade ist; oder durch
\[
\pm\frac{1\mathbin{.}3\mathbin{.}5\mathinner{.\,.\,.\,.} n}{(n+2)(n+4)\mathinner{.\,.\,.\,.}(2n+1)}F(-\tfrac12(n+1),\ \tfrac12n+1,\ \tfrac12,\ uu)
\]
für ungerades $n$, wobei das obere oder untere Zeichen gilt, je nachdem, ob $\frac12(n+1)$ gerade oder ungerade ist.

Die Funktion $U'$ lässt keinen ebenso einfachen allgemeinen Ausdruck zu. Aus der Entstehung der Größen $V,V',V''$ usw. selbst ergibt sich jedoch leicht, dass das letzte Glied von $U'$ für gerades $n$ lautet
\[
=\pm\frac{2\mathbin{.}2\mathbin{.}4\mathbin{.}4\mathbin{.}6\mathbin{.}6\mathinner{.\,.\,.\,.\,.} n\mathbin{.}n}{3\mathbin{.}5\mathbin{.}7\mathbin{.}9\mathbin{.}11\mathbin{.}13\mathinner{.\,.\,.\,.\,.}(2n-1)(2n+1)}
\]
wobei das obere oder untere Zeichen gilt, je nachdem, ob $\frac12n$ gerade oder ungerade ist.

Die Funktion $U''=\varphi W^{(n+1)}-V^{(n+1)}$, deren erstes Glied wir bereits im vorigen Artikel zu bestimmen gelehrt haben, lässt sich auch durch einen Rekursionsalgorithmus entwickeln, da offenbar allgemein gilt:
\[
\begin{aligned}
\varphi W''-V''&=w'(\varphi W'-V')+v'(\varphi W-V)\\
\varphi W'''-V'''&=w''(\varphi W''-V'')+v''(\varphi W'-V')\\
\varphi W''''-V''''&=w'''(\varphi W'''-V''')+v'''(\varphi W''-V'')
\end{aligned}
\]
\SourceLeaf{198}
usw., und somit in dem von uns betrachteten Fall
\[
\begin{aligned}&\varphi W^{(m+2)}-V^{(m+2)}=u(\varphi W^{(m+1)}-V^{(m+1)})\\[5pt]&-\frac{(m+1)^2}{(2m-1)(2m+1)}(\varphi W^{(m)}-V^{(m)})\end{aligned}
\]
So finden wir
\[
\begin{aligned}
\varphi W-V&=u^{-1}+\tfrac13u^{-3}+\tfrac15u^{-5}+\tfrac17u^{-7}+\text{usw.}\\
\varphi W'-V'&=\tfrac13u^{-2}+\tfrac15u^{-4}+\tfrac17u^{-6}+\tfrac19u^{-8}+\text{usw.}\\
\varphi W''-V''&=\tfrac4{45}u^{-3}+\tfrac8{105}u^{-5}+\tfrac4{63}u^{-7}+\tfrac{112}{2079}u^{-9}+\text{usw.}\\
\varphi W'''-V'''&=\tfrac4{175}u^{-4}+\tfrac8{315}u^{-6}+\tfrac4{165}u^{-8}+\tfrac{16}{715}u^{-10}+\text{usw.}
\end{aligned}
\]
usw. Diese Reihen lassen sich auch folgendermaßen darstellen:
\[
\begin{aligned}
\varphi W-V&=u^{-1}\left(1+\frac{1\mathbin{.}2}{2\mathbin{.}3}u^{-2}\right.\\[3pt]&\qquad\left.+\frac{1\mathbin{.}2\mathbin{.}3\mathbin{.}4}{2\mathbin{.}4\mathbin{.}3\mathbin{.}5}u^{-4}+\frac{1\mathbin{.}2\mathbin{.}3\mathbin{.}4\mathbin{.}5\mathbin{.}6}{2\mathbin{.}4\mathbin{.}6\mathbin{.}3\mathbin{.}5\mathbin{.}7}u^{-6}+\text{usw.}\right)\\[4pt]
\varphi W'-V'&=\tfrac13u^{-2}\left(1+\frac{2\mathbin{.}3}{2\mathbin{.}5}u^{-4}\right.\\[3pt]&\qquad\left.+\frac{2\mathbin{.}3\mathbin{.}4\mathbin{.}5}{2\mathbin{.}4\mathbin{.}5\mathbin{.}7}u^{-4}+\frac{2\mathbin{.}3\mathbin{.}4\mathbin{.}5\mathbin{.}6\mathbin{.}7}{2\mathbin{.}4\mathbin{.}6\mathbin{.}5\mathbin{.}7\mathbin{.}9}u^{-6}+\text{usw.}\right)\\[4pt]
\varphi W''-V''&=\tfrac4{45}u^{-3}\left(1+\frac{3\mathbin{.}4}{2\mathbin{.}7}u^{-2}\right.\\[3pt]&\qquad\left.+\frac{3\mathbin{.}4\mathbin{.}5\mathbin{.}6}{2\mathbin{.}4\mathbin{.}7\mathbin{.}9}u^{-4}+\frac{3\mathbin{.}4\mathbin{.}5\mathbin{.}6\mathbin{.}7\mathbin{.}8}{2\mathbin{.}4\mathbin{.}6\mathbin{.}7\mathbin{.}9\mathbin{.}11}u^{-6}+\text{usw.}\right)\\[4pt]
\varphi W'''-V'''&=\tfrac4{175}u^{-4}\left(1+\frac{4\mathbin{.}5}{2\mathbin{.}9}u^{-2}\right.\\[3pt]&\qquad\left.+\frac{4\mathbin{.}5\mathbin{.}6\mathbin{.}7}{2\mathbin{.}4\mathbin{.}9\mathbin{.}11}u^{-4}+\frac{4\mathbin{.}5\mathbin{.}6\mathbin{.}7\mathbin{.}8\mathbin{.}9}{2\mathbin{.}4\mathbin{.}6\mathbin{.}9\mathbin{.}11\mathbin{.}13}u^{-6}+\text{usw.}\right)
\end{aligned}
\]
usw. Folgen wir dieser Induktion, so ist allgemein
\[
U''=\varphi W^{(n+1)}-V^{(n+1)}\quad\text{gleich dem Produkt aus}
\]
\[
\frac{2\mathbin{.}2\mathbin{.}3\mathbin{.}3\mathbin{.}4\mathbin{.}4\mathinner{.\,.\,.\,.\,.}(n+1)\mathbin{.}(n+1)}{3\mathbin{.}3\mathbin{.}5\mathbin{.}5\mathbin{.}7\mathbin{.}7\mathbin{.}9\mathinner{.\,.\,.\,.\,.}(2n+1)(2n+3)}u^{-(n+2)}
\]
und der unendlichen Reihe
\[
1+\frac{(n+2)(n+3)}{2(2n+5)}u^{-2}+\frac{(n+2)(n+3)(n+4)(n+5)}{2\mathbin{.}4\mathbin{.}(2n+5)(2n+7)}u^{-4}+\text{usw.}
\]
oder, wenn man es vorzieht, der Funktion $F(\frac12n+1,\ \frac12n+\frac32,\ n+\frac52,\ u^{-2})$. Auch diese Induktion lässt sich sehr leicht zur vollen Gewissheit erheben, entweder durch die allgemein bekannte Methode oder mithilfe der Formel 19 in der mehrfach angeführten Abhandlung.

